% Preamble for Mechanics Vol 1 (loaded by main.tex)

% ============================================================================
%  MECHANICS — Volume I
%  Formal book design.  Monochrome text, rule-based structure, colour
%  reserved entirely for the illustrations.
% ============================================================================

\usepackage[fontsize=10.5pt]{fontsize}  % body 10.5pt: a formal A4 textbook page
\usepackage[T1]{fontenc}
\usepackage[utf8]{inputenc}
\usepackage{mathpazo}                 % Palatino text + matching mathematics
\usepackage[scaled=0.88]{helvet}      % sans, used only where unavoidable
\linespread{1.0}

\usepackage[a4paper,twoside,inner=25mm,outer=22mm,top=22mm,bottom=25mm,
            headsep=6mm,footskip=12mm]{geometry}

\usepackage{amsmath,amssymb,bm,mathtools}
\usepackage{siunitx}
% units for Volume III (eV, keV, MeV, Bq, Gy, Sv come with siunitx)
\DeclareSIUnit{\amu}{u}              % atomic mass unit
\DeclareSIUnit{\curie}{Ci}
\DeclareSIUnit{\year}{y}
\DeclareSIUnit{\fm}{fm}                % femtometre (not built in)
\usepackage{tikz}
\usepackage{pgfplots}
\pgfplotsset{compat=1.18}
\usetikzlibrary{arrows.meta,decorations.pathmorphing,decorations.markings,%
  decorations.pathreplacing,angles,quotes,patterns,patterns.meta,calc,positioning,%
  shapes.geometric,shadings,fadings,backgrounds,shapes.misc,hobby,bending,3d,perspective}
\usepgfplotslibrary{fillbetween}
\usepackage[most]{tcolorbox}
\usepackage{xcolor}
\usepackage{multicol}
\usepackage{enumitem}
\usepackage{graphicx}
\usepackage{cancel}
\usepackage[explicit]{titlesec}
\usepackage{titletoc}
\usepackage{fancyhdr}
\usepackage{booktabs}
\usepackage{needspace}
\usepackage{flafter}
\usepackage[section]{placeins}
\usepackage{emptypage}
\usepackage[expansion=false]{microtype}
\usepackage{hyperref}

% ------------------------------------------------------- native figures ---
% Figures drawn in TikZ/PGFPlots live one per file in figures/<part>/ (part =
% foundations, vol1, vol2, vol3) and share the palette and styles of
% figures/bookstyle.tex.  Include one by name alone (\tikzfig finds the folder):
%   \tikzfig{grav_shell}
% inside a figure float.  With shell-escape (latexmk does this via .latexmkrc)
% each picture is compiled once and cached in figures-cache/; without it the
% pictures are simply drawn inline, so the book always builds.
% =============================================================================
%  bookstyle.tex: the one visual vocabulary shared by every figure in the book
% =============================================================================
%  Loaded by preamble.tex (and by preview.tex for drawing a figure on its own).
%  Every TikZ/PGFPlots figure in figures/ draws with these names only:
%  no raw HTML colours, no ad hoc line widths, no font sizes set by hand.
%
%  Rules of thumb
%   * Draw at the printed size: never scale a figure with width= or scale=
%     after the fact, so text stays at the sizes set here.  Labels are \small
%     (9.6pt with the 10.5pt body), which puts their subscripts at 6.7pt, above the
%     book's 6.5pt floor;
%     \footnotesize is the floor for plain text and must not carry subscripts.
%   * Colour encodes meaning, and the meaning is the same in every chapter
%     (the book-wide legend; a figure may depart from it only to pair or
%     distinguish several vectors of one kind, e.g. third-law pairs or Lami's
%     three forces, and then the caption says so):
%       crimson  every force: weight, normal reaction, tension, gravity, F
%       accent   friction and applied pushes/pulls, where they must be told
%                apart from the contact forces; kinetic-energy bars (K)
%       forest   velocity v and angular velocity omega
%       primary  acceleration a and angular acceleration alpha; main plot curve
%       teal     momentum p and angular momentum L; components and legs
%       plum     potential energy U, total energy E, potentials
%       steel    dimensions, construction lines, secondary labels
%       grey     guides, paths not taken, ghost positions
%     Thermal physics (Volume II):
%       accent   heat Q, heat current H; isochoric curves (an isochoric step
%                exchanges only heat)
%       plum     work W
%       primary  isotherms (crimson/primary for hot/cold isotherms in a cycle)
%       teal     adiabats;  forest  isobars
%       reservoirs: hot cAmber!55 fill, cold cEarth fill
%     Nuclear physics (Volume III):
%       proton   pale accent disc (style `proton')   neutron  cMoon disc (`neutron')
%       alpha    plum      beta-minus, electron  primary      beta-plus, positron  crimson
%       neutrino teal, dashed      gamma, photon  accent, always a wavy line (`photon')
%   * Labels sit beside what they name, in the colour of what they name.
% -----------------------------------------------------------------------------

% ---------------------------------------------------------------- palette ---
% Prefixed with c so they never collide with the text greys in preamble.tex
% (where `ink' is pure black for the prose).
\definecolor{cInk}{HTML}{102A43}     % outlines, axes, neutral labels
\definecolor{cPrim}{HTML}{1F6FEB}    % blue
\definecolor{cTeal}{HTML}{0E7C7B}
\definecolor{cAcc}{HTML}{E85D04}     % orange
\definecolor{cCrim}{HTML}{C1121F}    % red
\definecolor{cFor}{HTML}{2D6A4F}     % green
\definecolor{cSteel}{HTML}{486581}
\definecolor{cPlum}{HTML}{7B2CBF}
\definecolor{cAmber}{HTML}{F4A261}
\definecolor{cGrey}{HTML}{9FB3C8}
% fills (pale, so outlines and arrows read on top of them)
\definecolor{cEarth}{HTML}{BFD3E6}   % planets
\definecolor{cMoon}{HTML}{D9E2EC}    % small bodies, light blocks
\definecolor{cBlock}{HTML}{93C5C6}   % blocks, bodies under study
\definecolor{cSlope}{HTML}{E1E7EF}   % inclines, ground fill
\definecolor{cRock}{HTML}{9AA8B6}    % mountains, walls, fixed supports
\definecolor{cSun}{HTML}{F4A261}

% ------------------------------------------------------------------ lines ---
\tikzset{
  % line weights, used everywhere instead of numbers
  outline/.style={line width=0.8pt,draw=cInk},
  thin line/.style={line width=0.5pt},
  hair line/.style={line width=0.35pt},
  curve/.style={line width=1.3pt},
  % ------------------------------------------------------------- arrows -----
  % one arrow head for the whole book; size scales with the line it ends
  bookhead/.tip={Stealth[length=2.4mm,width=1.7mm]},
  smallhead/.tip={Stealth[length=1.8mm,width=1.3mm]},
  vector/.style={-bookhead,line width=1.1pt,line cap=round},
  force/.style={vector,cCrim},
  applied/.style={vector,cAcc},
  velocity/.style={vector,cFor},
  accel/.style={vector,cPrim},
  momentum/.style={vector,cTeal},
  % dimension lines and leaders
  dim/.style={{smallhead}-{smallhead},thin line,cSteel},
  dimtick/.style={thin line,cSteel,densely dashed},
  leader/.style={hair line},
  guide/.style={thin line,cGrey,dash pattern=on 3pt off 2pt},
  orbitpath/.style={thin line,cGrey,dash pattern=on 4pt off 3pt},
  % --------------------------------------------------------------- bodies ---
  body/.style={fill=cBlock,draw=cSteel,line width=0.6pt},
  planet/.style={fill=cEarth,draw=cInk,line width=0.9pt},
  moon/.style={fill=cMoon,draw=cInk,line width=0.7pt},
  point mass/.style={circle,fill=cInk,inner sep=0pt,minimum size=3.2pt},
  % `ground' hatches a CLOSED region (a thin rectangle under the surface);
  % on a bare line it draws only the line.  For the usual pale strip under a
  % surface use `ground strip' on a rectangle, then draw the surface with `outline'.
  ground/.style={line width=0.8pt,draw=cInk,
    postaction={draw=none,pattern={Lines[angle=45,distance=2.2pt,line width=0.35pt]},
      pattern color=cInk}},
  ground strip/.style={fill=cSlope,draw=none},
  % a thin curved arrow for omega, alpha and other senses of rotation
  rotation/.style={-smallhead,thin line,cInk},
  % a small square marking a right angle; use as  \draw[right angle] (corner) ++(...)
  right angle/.style={hair line,cInk},
  % a medium-weight construction line in any colour: [construction,cFor]
  construction/.style={line width=0.8pt},
  % a dot in any colour: \node[dot,cCrim] at (P) {};
  dot/.style={circle,fill,inner sep=0pt,minimum size=3.2pt},
  % angle arcs: \draw[angle arc] (A) arc[...];  label it with a plain node
  angle arc/.style={thin line,cInk},
  % Riemann strips and shaded areas: pale fill, darker hairline edge
  strip/.style={fill=cPrim!14,draw=cPrim!60,hair line},
  area fill/.style={fill=cPrim!12,draw=none},
  % components of a vector (dashed, same colour as the vector, lighter head)
  component/.style={-smallhead,thin line,dash pattern=on 3pt off 1.5pt},
  % a resultant built by parallelogram: dashed, full weight
  resultant/.style={dash pattern=on 4pt off 2pt},
  parallelogram/.style={hair line,cGrey,dash pattern=on 2pt off 1.5pt},
  % rotation axes and hidden (back) edges of solids
  rotation axis/.style={thin line,cInk,dash pattern=on 6pt off 2pt on 1pt off 2pt},
  hidden/.style={dash pattern=on 2pt off 1.5pt},
  % the slice in a derivation figure (strip, ring, disc, shell)
  slice/.style={fill=cAcc!40,draw=cAcc,thin line},
  % a ghost: an earlier or later position of the same body.  Opacity also fades
  % any node placed on the same path, so put a ghost's label in a separate node.
  ghost/.style={opacity=0.35},
  % a particle or bob: shaded ball with a steel rim;  \pic at (P) {particle};
  particle/.pic={\shade[ball color=cBlock] (0,0) circle[radius=0.11];
                 \draw[cSteel,thin line] (0,0) circle[radius=0.11];},
  % stable and unstable equilibrium points on a U(x) curve
  stable point/.style={circle,fill=cInk,inner sep=0pt,minimum size=4pt},
  unstable point/.style={circle,draw=cInk,fill=white,thin line,inner sep=0pt,minimum size=4pt},
  % energy bars
  U bar/.style={fill=cPlum!55,draw=cPlum,thin line},
  K bar/.style={fill=cAcc!55,draw=cAcc,thin line},
  % horizontal circles drawn in perspective use this y-squash for the ellipse
  % (x radius r, y radius \persp*r)
  % the centre-of-mass marker: \pic at (G) {cm};
  cm/.pic={\fill[white] (0,0) circle[radius=0.09];
           \fill[cInk] (0,0) -- (0.09,0) arc[start angle=0,end angle=90,radius=0.09] -- cycle;
           \fill[cInk] (0,0) -- (-0.09,0) arc[start angle=180,end angle=270,radius=0.09] -- cycle;
           \draw[cInk,hair line] (0,0) circle[radius=0.09];},
  % out of / into the page
  out of page/.pic={\draw[cInk,thin line,fill=white] (0,0) circle[radius=0.09];
                    \fill[cInk] (0,0) circle[radius=0.025];},
  into page/.pic={\draw[cInk,thin line,fill=white] (0,0) circle[radius=0.09];
                  \draw[cInk,thin line] (-0.064,-0.064) -- (0.064,0.064) (-0.064,0.064) -- (0.064,-0.064);},
  % caution: `knockout' fills with the CURRENT colour if a bare colour name is
  % given alongside it; colour knocked-out text with text=<colour>.
  % the contact point of a rolling body
  contact point/.style={circle,fill=cInk,inner sep=0pt,minimum size=2.6pt},
  % an arrowhead part-way along a path (cycle directions, flow lines):
  % \draw[curve,mid arrow] ... or [mid arrow=0.56]
  mid arrow/.style={postaction={decorate},decoration={markings,
    mark=at position #1 with {\arrow{bookhead}}}},
  mid arrow/.default=0.5,
  % ---------------------------------------------------- nuclear particles ---
  % nucleons are pale discs (size set by the figure: circle,proton,minimum size=..);
  % moving particles are arrows or outlined dots; a photon is always a wavy line,
  % so it is never confused with a proton disc although both use the accent
  proton/.style={draw=cAcc,fill=cAcc!35,thin line},
  neutron/.style={draw=cSteel,fill=cMoon,thin line},
  photon/.style={cAcc,decorate,decoration={snake,amplitude=1.3pt,
    segment length=5pt,post length=2.5mm}},
  % ---------------------------------------------------------------- text ----
  % every figure starts from \small (9.6pt; subscripts 6.7pt); \footnotesize (8.75pt)
  % is the floor, for plain words only
  figure text/.style={font=\small,text=cInk},
  small label/.style={font=\footnotesize},
  note/.style={font=\small\itshape,text=cInk},
  % a label that must sit over a line: white knock-out behind the text
  knockout/.style={fill=white,inner sep=1.2pt},
  % applied to every figure by \tikzfig
  % default drawing colour is ink, so plain `dot', `fill' and `draw' match the palette
  bookfigure/.style={figure text,color=cInk,line join=round,>=bookhead},
}

\def\persp{0.25}
% ----------------------------------------------------------------- plots ----
\usepgfplotslibrary{groupplots}
\pgfplotsset{
  % the default look of every plot: two axis lines with arrow tips, no box,
  % no grid; ticks only where the text refers to a value
  bookplot/.style={
    axis lines=left,
    axis line style={line width=0.7pt,draw=cInk,-{Stealth[length=2mm,width=1.5mm]}},
    every tick/.style={line width=0.5pt,draw=cInk},
    tick label style={font=\small,text=cInk},
    label style={font=\small,text=cInk},
    every axis plot/.append style={curve},
    xtick=\empty, ytick=\empty,
    clip=false,
    samples=120,
    enlargelimits=false,
    axis on top=false,
  },
  % a zero line drawn inside a plot (energy plots)
  zero line/.style={thin line,cInk},
}

\usetikzlibrary{external}
\newif\iffigcache
\ifdefined\pdfshellescape\ifnum\pdfshellescape=1 \figcachetrue\fi\fi
\iffigcache
  \tikzexternalize[prefix=figures-cache/]
  % cache only \tikzfig pictures, never the tcolorbox frames or inline marks
  \AtBeginDocument{\tikzexternaldisable}
\fi
\newcommand{\tikzfig}[1]{%
  \begingroup
  \iffigcache\tikzexternalenable\tikzsetnextfilename{#1}\fi
  \tikzset{every picture/.append style={bookfigure}}%
  \figfile{#1}\input{\thefigfile}\endgroup}
% the folder of a figure: the first of figures/{foundations,vol1,vol2,vol3}/ holding NAME.tex
\makeatletter
\newcommand{\figfile}[1]{\def\thefigfile{figures/#1}%
  \@for\figdir:=foundations,vol1,vol2,vol3\do{%
    \IfFileExists{figures/\figdir/#1.tex}{\edef\thefigfile{figures/\figdir/#1}}{}}}
\makeatother

% ---------------------------------------------------------------- palette ---
% Text is black.  These two greys are for rules and secondary marks only;
% every saturated colour in this book lives inside an illustration.
\definecolor{ink}{gray}{0.00}
\definecolor{rulegrey}{gray}{0.35}
\definecolor{lightrule}{gray}{0.72}
\definecolor{captiongrey}{gray}{0.25}

% legacy colour names kept so that existing figure/TikZ code still compiles;
% all are mapped onto the monochrome scheme for anything that sets type.
\definecolor{inkblue}{gray}{0.00}
\definecolor{primary}{gray}{0.00}
\definecolor{accent}{gray}{0.00}
\definecolor{teal}{gray}{0.00}
\definecolor{plum}{gray}{0.00}
\definecolor{crimson}{gray}{0.00}
\definecolor{forest}{gray}{0.00}
\definecolor{amber}{gray}{0.45}
\definecolor{softgrey}{gray}{1.00}
\definecolor{rule1}{gray}{0.72}

\hypersetup{colorlinks=true,linkcolor=black,urlcolor=black,citecolor=black,
            bookmarksnumbered=true,pdfborder={0 0 0},
            pdftitle={Physics: A Reasoning-First Companion},
            pdfsubject={Physics for JEE and NEET: mechanics, properties of matter, thermal physics, nuclei and radioactivity},
            pdfkeywords={physics, mechanics, fluids, thermodynamics, kinetic theory, JEE, NEET, Class 11}}

\tikzset{
  lbl/.style={fill=white,draw=none,inner sep=2pt,rounded corners=1.5pt,font=\footnotesize},
  lblbox/.style={fill=white,draw=lightrule,line width=0.4pt,inner sep=2.5pt,rounded corners=1.5pt,font=\footnotesize},
  meas/.style={rulegrey,line width=0.7pt},
  dimchip/.style={fill=white,draw=lightrule,line width=0.4pt,inner sep=2.2pt,rounded corners=1.5pt,font=\footnotesize\bfseries},
  vec/.style={-{Stealth[length=3mm,width=2.2mm]},line width=1.1pt},
  thinvec/.style={-{Stealth[length=2.4mm]},line width=0.8pt},
  dashline/.style={dashed,line width=0.6pt,gray},
  axis line/.style={-{Stealth[length=2.5mm]},line width=0.8pt,black!70}
}
\newcommand{\circled}[1]{\tikz[baseline=(c.base)]{\node[shape=circle,draw=black,
  line width=0.5pt,inner sep=1pt,minimum size=1.4ex,font=\scriptsize\bfseries](c){#1};}}

\raggedbottom
% no single lines stranded at the top or bottom of a page, and no page
% that ends on a hyphen
\widowpenalty=10000
\clubpenalty=10000
\displaywidowpenalty=10000
\brokenpenalty=10000
\setlength{\textfloatsep}{10pt plus 3pt minus 2pt}
\setlength{\floatsep}{8pt plus 2pt minus 2pt}
\setlength{\intextsep}{8pt plus 2pt minus 2pt}
\renewcommand{\topfraction}{0.85}
\renewcommand{\bottomfraction}{0.6}
\renewcommand{\textfraction}{0.12}
\renewcommand{\floatpagefraction}{0.72}
\setlength{\parskip}{0pt}
\setlength{\parindent}{1.2em}
\setlist[itemize]{leftmargin=1.4em,topsep=4pt,itemsep=2.5pt,parsep=0pt}
\setlist[enumerate]{leftmargin=1.6em,topsep=4pt,itemsep=3.5pt,parsep=0pt}

% ======================================================= sectional units ====
\renewcommand{\thepart}{\Alph{part}}

\titleformat{\chapter}[display]
  {\normalfont\filright}
  {\vspace*{6pt}\normalsize\scshape\chapterword\ \Huge\bfseries\thechapter}
  {6pt}
  {\vspace{2pt}{\color{rulegrey}\titlerule[1.1pt]}\vspace{10pt}%
   \parbox{\textwidth}{\raggedright\huge\bfseries\linespread{1.0}\selectfont#1}\vspace{6pt}%
   {\color{rulegrey}\titlerule[0.5pt]}}
\titlespacing*{\chapter}{0pt}{-12pt}{26pt}

\titleformat{\section}{\normalfont\Large\bfseries}{\thesection}{0.8em}{#1}
\titlespacing*{\section}{0pt}{14pt plus 3pt minus 2pt}{6pt}

\titleformat{\subsection}{\normalfont\large\bfseries}{\thesubsection}{0.7em}{#1}
\titlespacing*{\subsection}{0pt}{10pt plus 2pt minus 2pt}{4pt}

\titleformat{\subsubsection}{\normalfont\normalsize\bfseries\itshape}{}{0pt}{#1}
\titlespacing*{\subsubsection}{0pt}{11pt plus 2pt}{3pt}

% running heads
\pagestyle{fancy}\fancyhf{}
\renewcommand{\headrulewidth}{0.4pt}
\renewcommand{\footrulewidth}{0pt}
\fancyhead[LE]{\small\scshape\nouppercase{\leftmark}}
\fancyhead[RO]{\small\scshape\nouppercase{\rightmark}}
\fancyhead[RE,LO]{\small\thepage}
\renewcommand{\chaptermark}[1]{\markboth{\thechapter.\ #1}{}}
\renewcommand{\sectionmark}[1]{\markright{\thesection\ \ #1}}
\fancypagestyle{plain}{\fancyhf{}\renewcommand{\headrulewidth}{0pt}%
  \fancyfoot[C]{\small\thepage}}

% ----------------------------------------------------------------- TOC ------
\setcounter{tocdepth}{1}
\setcounter{secnumdepth}{3}
% \volumetocline{numeral}{title}{page}: the top level of the contents
\newcommand{\volumetocline}[3]{\par\addvspace{20pt}\noindent
  {\Large\scshape Volume #1\quad\bfseries\upshape #2\hfill\normalsize\mdseries #3}\par\addvspace{2pt}}
\titlecontents{part}[0pt]
  {\addvspace{14pt}\large\bfseries}{}{}{}[\addvspace{4pt}]
\titlecontents{chapter}[1.6em]
  {\addvspace{6pt}\bfseries}
  {\contentslabel{1.6em}}{}
  {\titlerule*[0.6pc]{.}\contentspage}[\addvspace{2pt}]
\titlecontents{section}[4.2em]
  {\small}
  {\contentslabel{2.9em}}{}
  {\titlerule*[0.5pc]{.}\contentspage}

% =========================================================== environments ===
% One visual grammar for the whole book:
%   * a statement that must be remembered       -> rules above and below
%   * a derivation                              -> hairline rule at the left
%   * an example, a caution, an aside           -> small type, hairline rule
% Nothing is filled with colour; weight and white space carry the hierarchy.

\newcounter{bxcount}[chapter]
\renewcommand{\thebxcount}{\thechapter.\arabic{bxcount}}

% --- rules above and below: definitions and key results ---------------------
\newtcolorbox{statementbox}[2]{enhanced,breakable,colback=white,colframe=white,
  boxrule=0pt,sharp corners,left=0pt,right=0pt,top=7pt,bottom=7pt,
  before skip=13pt,after skip=7pt,
  borderline north={0.9pt}{0pt}{ink},
  borderline south={0.9pt}{0pt}{ink},
  overlay unbroken and first={%
    \node[anchor=north west,inner sep=0pt,font=\footnotesize\scshape,
          fill=white,inner xsep=3pt]
      at ([xshift=-3pt,yshift=5.5pt]frame.north west) {#1\ \ \textnormal{\textbf{#2}}};}}

\newenvironment{defbox}[1]
  {\refstepcounter{bxcount}%
   \begin{statementbox}{Definition \thebxcount}{#1}}
  {\end{statementbox}}

\newenvironment{keybox}[1]
  {\refstepcounter{bxcount}%
   \begin{statementbox}{Result \thebxcount}{#1}}
  {\end{statementbox}}

% --- hairline at the left: derivations --------------------------------------
\newtcolorbox{sidebox}[2]{enhanced,breakable,colback=white,colframe=white,
  boxrule=0pt,sharp corners,left=11pt,right=0pt,top=4pt,bottom=4pt,
  before skip=6.5pt,after skip=6.5pt,
  borderline west={0.8pt}{3pt}{ink},
  % The label is real text, not an overlay: an overlay's clearance (a \vspace)
  % is discarded when the box starts at the top of a page, and the body then
  % printed over the label.
  before upper={{\noindent\footnotesize\scshape #1\ \ \textnormal{\textbf{#2}}\par}%
    \nobreak\vspace{4pt}\nobreak}}

\newenvironment{derivbox}[1]
  {\refstepcounter{bxcount}%
   \begin{sidebox}{Derivation \thebxcount}{#1}}
  {\end{sidebox}}

\newenvironment{exbox}[1]
  {\refstepcounter{bxcount}%
   \begin{sidebox}{Example \thebxcount}{#1}\small}
  {\hfill$\square$\end{sidebox}}

\newenvironment{archbox}[1]
  {\refstepcounter{bxcount}%
   \begin{sidebox}{Problem type \thebxcount}{#1}\small}
  {\end{sidebox}}

\newenvironment{challengebox}[1]
  {\refstepcounter{bxcount}%
   \begin{sidebox}{Challenge \thebxcount}{#1}\small}
  {\end{sidebox}}

\newenvironment{methodbox}[1]
  {\refstepcounter{bxcount}%
   \begin{statementbox}{Method \thebxcount}{#1}}
  {\end{statementbox}}

\newenvironment{tablebox}[1]
  {\refstepcounter{bxcount}%
   \begin{statementbox}{Table \thebxcount}{#1}\small}
  {\end{statementbox}}

\newenvironment{remarkbox}[1]
  {\refstepcounter{bxcount}%
   \begin{sidebox}{Remark \thebxcount}{#1}\small}
  {\end{sidebox}}

% --- small type, light hairline: cautions, asides, motivation ---------------
\newtcolorbox{notebox}[1]{enhanced,breakable,colback=white,colframe=white,
  boxrule=0pt,sharp corners,left=10pt,right=0pt,top=2pt,bottom=2pt,
  before skip=6pt,after skip=6pt,
  borderline west={0.5pt}{3pt}{rulegrey},
  before upper={\small\textit{#1}\ \ignorespaces}}

\newenvironment{trapbox}[1]{\begin{notebox}{Caution: #1.}}{\end{notebox}}
\newenvironment{motivbox}[1]{\begin{notebox}{#1.}}{\end{notebox}}
\newenvironment{summarybox}{\begin{notebox}{What to carry forward.}}{\end{notebox}}

\newcommand{\why}[1]{\begin{notebox}{Why this matters.}#1\end{notebox}}
\newcommand{\inwords}[1]{\begin{notebox}{In words.}#1\end{notebox}}
\newcommand{\scopebox}[1]{\begin{notebox}{Syllabus note.}#1\end{notebox}}

% --- figures ----------------------------------------------------------------
% Inline, unframed, centred.  The caption carries the number.
\newenvironment{figbox}
  {\par\addvspace{12pt}\begingroup\centering}
  {\par\endgroup\addvspace{3pt}}

\newcounter{figc}[chapter]
\renewcommand{\thefigc}{\thechapter.\arabic{figc}}
\newcommand{\figcaption}[2]{%
  \refstepcounter{figc}\label{#1}%
  \par\vspace{3pt}%
  {\centering\begin{minipage}{0.93\textwidth}%
   \footnotesize\color{captiongrey}\textbf{Figure~\thefigc.}\ #2\par
   \end{minipage}\par}%
  \vspace{7pt}}
\newcommand{\figcap}[1]{\figcaption{autofig\thefigc}{#1}}

% --- run-in marks -----------------------------------------------------------
\newcommand{\stepn}[1]{\par\medskip\noindent\textbf{Step #1.}\ \ignorespaces}
% \ans{...}: the answer sits flush right on the last line if it fits there,
% otherwise flush right on a line of its own (TeXbook, p. 106)
\newcommand{\ans}[1]{\ifvmode\par\noindent\fi\unskip\nobreak\hfil\penalty50
  \hskip1.5em\hbox{}\nobreak\hfill\textbf{[#1]}{\parfillskip=0pt\finalhyphendemerits=0\par}}
\newcommand{\archans}[2]{\gdef\archanswerstore{\par\smallskip\noindent\textbf{Answer: #1.}\ #2}}
\let\archanswerstore\relax
\newcommand{\challengeans}[1]{\par\medskip\noindent\textbf{Answer.}\ #1}

\newcommand{\arch}[6]{\begin{archbox}{#1}%
\textbf{Tell.}\ #2\par\textbf{Method.}\ #3\par\textbf{Trap.}\ #4\par
\textbf{Try it.}\ #5 \ans{#6}\end{archbox}}

\newcommand{\archq}[9]{\begin{archbox}{#1}%
\textbf{Tell.}\ #2\par\textbf{Method.}\ #3\par\textbf{Trap.}\ #4\par
\smallskip\textbf{Exam-style, single correct.}\ #5\par\smallskip
\begin{tabular}{@{}p{0.45\linewidth}@{\hspace{8pt}}p{0.45\linewidth}@{}}
(A)~#6 & (B)~#7\\[1pt](C)~#8 & (D)~#9
\end{tabular}\archanswerstore\global\let\archanswerstore\relax\end{archbox}}

\newcommand{\wex}[1]{\begin{exbox}{#1}}
\newcommand{\wexend}{\end{exbox}}

% --- volumes, chapter and part openers --------------------------------------
% The book is one PDF in several volumes.  The Foundations chapter (the
% mathematical toolkit) comes first and serves every volume; it is numbered F,
% so its boxes and figures read F.1, F.2.  Chapter and part numbers restart in
% each volume.  \currvol holds F, I, II, ...; hyperref's anchors include it,
% so Figure 2.3 of Volume I and Figure 2.3 of Volume II never collide.
\def\currvol{F}
\def\chapterword{chapter}

% \foundationschapter{Title}: the shared opening chapter
\newcommand{\foundationschapter}[1]{%
  \def\currvol{F}\def\chapterword{foundations}%
  \renewcommand{\thechapter}{F}%
  \setcounter{chapter}{0}%
  \chapter{#1}}

% \volumedivider{numeral}{Title}{thesis}{part list}: starts a volume
\newcommand{\volumedivider}[4]{%
  \cleardoublepage
  \def\currvol{#1}\def\chapterword{chapter}%
  \renewcommand{\thechapter}{\arabic{chapter}}%
  \setcounter{chapter}{0}\setcounter{part}{0}%
  \thispagestyle{empty}%
  \phantomsection
  \pdfbookmark[-2]{Volume #1: #2}{volume.#1}%
  \addtocontents{toc}{\protect\volumetocline{#1}{#2}{\thepage}}%
  \vspace*{4.2cm}
  \begin{center}
    {\Large\scshape Volume #1}\\[6mm]
    {\color{rulegrey}\rule{0.5\textwidth}{0.9pt}}\\[9mm]
    {\Huge\bfseries #2}\\[10mm]
    {\color{rulegrey}\rule{0.5\textwidth}{0.4pt}}\\[10mm]
    \begin{minipage}{0.76\textwidth}\centering\large\itshape #3\end{minipage}\\[14mm]
    \begin{minipage}{0.80\textwidth}\centering #4\end{minipage}
  \end{center}
  \cleardoublepage}

% \volumeopening{Title}: an unnumbered opening chapter of a volume (purpose,
% map, notation).  It carries no numbered boxes or figures.
\newcommand{\volumeopening}[1]{%
  \chapter*{#1}\addcontentsline{toc}{chapter}{#1}\markboth{#1}{#1}}

% \ncertchapter{n}{Title} sets the chapter number explicitly.
\newcommand{\ncertchapter}[2]{%
  \setcounter{chapter}{\numexpr#1-1\relax}%
  \chapter{#2}}

% \partdivider{Letter}{Title}{thesis}{chapter list}
\newcommand{\partdivider}[4]{%
  \cleardoublepage
  \refstepcounter{part}%
  \thispagestyle{empty}%
  \vspace*{5.5cm}
  \begin{center}
    {\normalsize\scshape Part \thepart}\\[5mm]
    {\color{rulegrey}\rule{0.42\textwidth}{0.9pt}}\\[7mm]
    {\huge\bfseries #2}\\[9mm]
    {\color{rulegrey}\rule{0.42\textwidth}{0.4pt}}\\[8mm]
    \begin{minipage}{0.76\textwidth}\centering\itshape #3\end{minipage}\\[11mm]
    \begin{minipage}{0.80\textwidth}\centering\small #4\end{minipage}
  \end{center}
  \addcontentsline{toc}{part}{\texorpdfstring{Part \thepart\quad #2}{Part \thepart: #2}}%
  \cleardoublepage}

% ---------------------------------------------------------------- macros ----
% cross-references to chapters: never type a chapter number (labels ch:toolkit,
% ch:measurement, ch:straightline, ch:vectors, ch:projectiles, ch:relative,
% ch:circkin, ch:newton, ch:connected, ch:friction, ch:circdyn, ch:energy,
% ch:collisions, ch:com, ch:rotation, ch:rolling, ch:gravitation, ch:kepler)
% Every chapter label is set with \chlabel, which also records its volume, so
% \chref prints "Chapter~11" within a volume, "Volume~I, Chapter~11" from
% another volume, and "the Foundations" for the toolkit.
% Volume II labels: ch:solids, ch:fluids, ch:thermal, ch:thermo, ch:kinetic.
% Volume III labels: ch:nucleus, ch:decay, ch:decaylaw, ch:reactions, ch:nucapps.
\makeatletter
\newcommand{\setchvol}[2]{\expandafter\gdef\csname chvol@#1\endcsname{#2}}
\newcommand{\chlabel}[1]{\label{#1}%
  \immediate\write\@auxout{\string\setchvol{#1}{\currvol}}}
\def\chv@F{F}
\newcommand{\chref}[1]{%
  \ifcsname chvol@#1\endcsname
    \expandafter\let\expandafter\chv@this\csname chvol@#1\endcsname
    \ifx\chv@this\chv@F the Foundations%
    \else\ifx\chv@this\currvol Chapter~\ref{#1}%
    \else Volume~\chv@this, Chapter~\ref{#1}\fi\fi
  \else Chapter~\ref{#1}\fi}
\makeatother
\makeatletter
\newcommand{\chsref}[2]{%
  \ifcsname chvol@#1\endcsname
    \expandafter\let\expandafter\chv@this\csname chvol@#1\endcsname
    \ifx\chv@this\currvol\else Volume~\chv@this, \fi
  \fi Chapters~\ref{#1} and~\ref{#2}}
\makeatother
\newcommand{\vc}[1]{\vec{\bm{#1}}}
\newcommand{\uhat}[1]{\hat{\bm{#1}}}
\newcommand{\A}{\vc{A}}
\newcommand{\B}{\vc{B}}
\newcommand{\R}{\vc{R}}
\newcommand{\ihat}{\uhat{\imath}}
\newcommand{\jhat}{\uhat{\jmath}}
\newcommand{\khat}{\uhat{k}}
\newcommand{\vv}{\vc{v}}
\newcommand{\va}{\vc{a}}
\newcommand{\vF}{\vc{F}}
\newcommand{\vp}{\vc{p}}
\newcommand{\vr}{\vc{r}}
\newcommand{\vL}{\vc{L}}
\newcommand{\vtau}{\vc{\tau}}
\newcommand{\vg}{\vc{g}}
\newcommand{\vomega}{\vc{\omega}}
\newcommand{\ex}{\times}
\newcommand{\faExclam}{}

% --- Volume III: nuclear notation ----------------------------------------------
% \nuc{A}{Z}{X} a nuclide with mass and atomic number; \iso{A}{X} mass number only
\newcommand{\nuc}[3]{\prescript{#1}{#2}{\mathrm{#3}}}
\newcommand{\iso}[2]{\prescript{#1}{}{\mathrm{#2}}}
\newcommand{\Act}{\mathcal{A}}      % activity (decays per second)
\newcommand{\Prod}{\mathcal{P}}     % constant production rate
\newcommand{\Thalf}{T_{1/2}}        % half-life
\newcommand{\Eb}{E_{\mathrm{b}}}    % binding energy

% --- unique hyperlink anchors across volumes ---------------------------------
\renewcommand{\theHpart}{\currvol.\arabic{part}}
\renewcommand{\theHchapter}{\currvol.\arabic{chapter}}
\renewcommand{\theHsection}{\theHchapter.\arabic{section}}
\renewcommand{\theHequation}{\theHchapter.\arabic{equation}}
\renewcommand{\theHfigure}{\theHchapter.\arabic{figure}}
\renewcommand{\theHtable}{\theHchapter.\arabic{table}}
\providecommand{\theHfigc}{}\renewcommand{\theHfigc}{\theHchapter.\arabic{figc}}
\providecommand{\theHbxcount}{}\renewcommand{\theHbxcount}{\theHchapter.\arabic{bxcount}}
